\documentclass[tikz]{standalone}
\usetikzlibrary{shapes}
\usetikzlibrary{math}


%% \foreach \n in {10,20,...,90}
%% do something with \n ;



\begin{document}
\begin{tikzpicture}

  %% Legs
  \tikzmath{
    \offset = .1; %% separation of legs
    \phase1=rand*100; %% phases of legs
    \phase2=rand*100; %% right to left, outside in.
    \phase3=rand*100;
    \phase4=rand*100;
    \phase5=rand*100;
    \phase6=rand*100;
    \phase7=rand*100;
    \phase8=rand*100;
  }
  \pgfmathsetseed{1}
  \draw[line width=10pt,line cap=round,black, domain=0:4, samples=100, variable=\x]
  plot ({.1*\x*sin (200*\x+\phase1)-7*\offset},{-\x+5});
  \draw[line width=8pt,line cap=round,orange!40!black, domain=0:4, samples=100, variable=\x]
  plot ({.025*\x^2*sin (200*\x+\phase1)-7*\offset},{-\x+5});

  \draw[line width=10pt,line cap=round,black, domain=0:4, samples=100, variable=\x]
  plot ({.025*\x^2*sin (200*\x+\phase2)+7*\offset},{-\x+5});
  \draw[line width=8pt,line cap=round,orange!80!yellow!40!black, domain=0:4, samples=100, variable=\x]
  plot ({.025*\x^2*sin (200*\x+\phase2)+7*\offset},{-\x+5});

  
  \draw[line width=10pt,line cap=round,black, domain=0:4, samples=100, variable=\x]
  plot ({.025*\x^2*sin (200*\x+\phase3)-5*\offset},{-\x+5});
  \draw[line width=8pt,line cap=round,orange!60!black, domain=0:4, samples=100, variable=\x]
  plot ({.025*\x^2*sin (200*\x+\phase3)-5*\offset},{-\x+5});

  \draw[line width=10pt,line cap=round,black, domain=0:4, samples=100, variable=\x]
  plot ({.025*\x^2*sin (200*\x+\phase4)+5*\offset},{-\x+5});
  \draw[line width=8pt,line cap=round,orange!80!yellow!60!black, domain=0:4, samples=100, variable=\x]
  plot ({.025*\x^2*sin (200*\x+\phase4)+5*\offset},{-\x+5});
  

  \draw[line width=10pt,line cap=round,black, domain=0:4, samples=100, variable=\x]
  plot ({.025*\x^2*sin (200*\x+\phase5)-3*\offset},{-\x+5});
  \draw[line width=8pt,line cap=round,orange!80!black, domain=0:4, samples=100, variable=\x]
  plot ({.025*\x^2*sin (200*\x+\phase5)-3*\offset},{-\x+5});

  \draw[line width=10pt,line cap=round,black, domain=0:4, samples=100, variable=\x]
  plot ({.025*\x^2*sin (200*\x+\phase6)+3*\offset},{-\x+5});
  \draw[line width=8pt,line cap=round,orange!80!yellow!80!black, domain=0:4, samples=100, variable=\x]
  plot ({.025*\x^2*sin (200*\x+\phase6)+3*\offset},{-\x+5});



  \draw[line width=10pt,line cap=round,black, domain=0:4, samples=100, variable=\x]
  plot ({.025*\x^2*sin (200*\x+\phase7)-\offset},{-\x+5});
  \draw[line width=8pt,line cap=round,orange, domain=0:4, samples=100, variable=\x]
  plot ({.025*\x^2*sin (200*\x+\phase7)-\offset},{-\x+5});

  \draw[line width=10pt,line cap=round,black, domain=0:4, samples=100, variable=\x]
  plot ({.025*\x^2*sin (200*\x+\phase8)+\offset},{-\x+5});
  \draw[line width=8pt,line cap=round,orange!80!yellow, domain=0:4, samples=100, variable=\x]
  plot ({.025*\x^2*sin (200*\x+\phase8)+\offset},{-\x+5});



  
  %% The main of the octolion is two stars
  \node[
    ultra thick,
    star,
    star points=11,
    star point height=.4cm,
    draw=black,
    fill=black!50!brown,
    minimum size=3cm,
  ] at (0,5.1) {};

  \node[
    ultra thick,
    rotate=180,
    star,
    star points=11,
    star point height=.5cm,
    draw=black,
    fill=yellow!30!brown,
    minimum size=3cm,
  ] at (0,5.1) {};


  %% Ears
  \draw[fill=orange,ultra thick](-.6,5.8) circle (.4);
  \draw[fill=black!50!brown,draw=none](-.6,5.8) circle (.25);

  \draw[fill=orange,ultra thick](.6,5.8) circle (.4);
  \draw[fill=black!50!brown,draw=none](.6,5.8) circle (.25);


  %% Head of lion
  \draw[fill=orange,draw=black,ultra thick] (0,5) ellipse (1.1cm and 1cm);


  %% eyes
  \draw[fill=black!50!blue,draw=black,ultra thick](-.6,5.1) circle (.2);
  \draw[fill=white!80!blue,draw=none](-.55,5.15) circle (.07);

  \draw[fill=black!50!green,draw=black,ultra thick](.6,5.1) circle (.2);
  \draw[fill=white!80!green,draw=none](.65,5.15) circle (.07);
  

  %% mouth/nose
  \draw[fill=white,draw=none,ultra thick] (0,4.5) ellipse (.58cm and .48cm);

  \draw [rounded corners,fill=black] (-.3,4.98)--(.3,4.98)--(0,4.6)--cycle;

  \draw[line width=2pt,rounded corners,cap=round] (0,4.7) --  (0,4.5) -- (-.22,4.3) -- (-.44,4.5);
  \draw[line width=2pt,rounded corners,cap=round] (0,4.7) --  (0,4.5) -- (.22,4.3) -- (.44,4.5);

 
  
  
\end{tikzpicture}
\end{document}
